% 从旧数学附录迁出的扫描、分治和约束说明；不新增算法。
\section{线段树维护单调栈：双向扫描与接雨水}\label{knowledge-seg-stack}
\texttt{MonotoneStackSeg} 维护区间加后的运行最大值贡献。
每段的 $F_0,F_1$ 分别为从左、从右扫描的最高值减原值之和。
若外部高度 $x$ 不低于整段最大值，贡献为 $x\cdot len-sum$。
否则，将孩子按扫描顺序记为 $A,B$：若 $x\ge\max A$，递归 $B$；
否则递归 $A$，而 $B$ 的贡献已经是 $F_{parent}-F_A$。
每次只递归一个孩子，因此单段带边界扫描为 $O(\log n)$。
区间平移不改变这些差值，懒标记只改最大值、原值和与标记本身。
区间修改回溯及任意子区间扫描最坏 $O(\log^2(n+1))$，建树 $O(n)$。

对全局前缀和后缀最大值 $L_i,R_i$，有 $\max(L_i,R_i)=M$，
其中 $M$ 为全局最大值。因此
\[
\sum_i\bigl(\min(L_i,R_i)-a_i\bigr)
=F_0+F_1-\bigl(nM-\sum_i a_i\bigr).
\]
根上即可 $O(1)$ 返回雨水体积，不必查询最大点位置。
整体平移不改变体积，负高度也成立；带外部高度的 scan 则明确使用调用者的边界。
所有乘法、和、标记及中间运算须在 long long 内，不能只检查最终答案。
附件原题为 NERC 2023 G，使用例属于比赛应用，不计正式模板题覆盖；
小规模有符号穷举、三十万次随机更新和二十万点规模本地验证通过，线上待补。

\section{动态 K-D Tree：矩形和与重建边界}\label{knowledge-kdt-sum}
\texttt{KDTreeSum} 按深度交替坐标轴，以另一轴打破平局。
重复坐标累加权值，其他点沿轴全序递归插入；孩子过重时收集编号并用
\texttt{nth\_element} 重建，复用原节点。
包围盒从本点初始化，仅合并非空孩子，不用有限正负无穷；
重建直接返回根编号，不用坐标 $-1$ 作为状态标志。
权值为零时仍保留节点，负权不影响几何剪枝。

每个孩子大小至多为父节点的 $3/4$，所以树高为 $O(\log(d+1))$；
矩形部分相交却可能递归两边，查询最坏仍为 $O(d)$。
按平均线性选择计，一次大小 $s$ 的重建需 $O(s\log(s+1))$，
插入摊还 $O(\log^2(d+1))$；不是每次操作的最坏对数界。
矩形使用闭边界，只比较坐标、不做坐标减法；
各点绝对权值和及更新中间量须在 long long 内。

P4148 的在线输入需逐参数 XOR 上次答案，驱动预留二十万个节点以避免扩容峰值。
本地默认 int 坐标节点为 56 字节，另有重建编号和运行环境开销；
该分配量不能代替线上 20 MB 内存通过记录。
独立 map 对拍、分割与包围盒审计、坐标极值及二十万规模已本地验证，线上待补。

\section{CDQ 卷积递推：贡献只跨越一次}\label{knowledge-cdq-conv}
\texttt{cdq\_convolution} 求
\[
f_i=b_i+\sum_{j<i}f_j g_{i-j},\qquad g_0=0.
\]
在半开区间 $[l,r)$ 中，先递归求完左半 $[l,m)$，
再将这段 $f$ 与 $g_{0\ldots r-l-1}$ 卷积，
把乘积的第 $i-l$ 项加入右半的 $f_i$，最后递归右半。
任意 $j<i$ 的依赖只有一个最小分治节点将二者分到左右两边，
所以每份贡献恰好加入一次。顺序不能改成先求两个孩子再卷积。

根层完整卷积结果长度为 $n+\lfloor n/2\rfloor-1$，
须不超过 NTT 的最大二次幂容量；单项输入不做变换，空输入返回空。
一般初始项 $b_i$ 均可使用，P4721 只取 $b_0=1$。
时间 $O(n\log^2(n+1))$，额外空间 $O(n)$，递归深度 $O(\log n)$。
附件标题为“分治 FFT”，实现实际依赖 NTT；
这里保留 CDQ 算法，不以形式幂级数求逆替代。
独立二次递推、多素数域、二十万项闭式参考及完整主程序已本地验证，线上待补。

\section{SPFA 费用流与残量方案}\label{knowledge-spfa-cost}
SPFA 备选每次在残量图上直接求最短路，保留负边、限流和继续增广，
但不允许负环。每次调用先将全部顶点以零距离入队检查，包括源点不可达的负环；
因为不连通的负环同样能降低总费用，不能忽略。
检查时不改流量，发现负环抛异常。记录严格改进路径的边数，达到顶点数即拒绝；
不要用松弛次数冒充路径边数。零费用环不发生严格改进，不会反复改写前驱。

沿最短残量路增广保持当前流量下的最小费用性。
反边可以撤销已经送过的流，最终方案应通过原边编号查询 used，不能累加历史路径当作最终流量。
flow 返回本次新增流量与费用，重复调用须同源汇；传零限流仍检查负环。
最小费用流允许负边与允许负环是两种不同的契约，后者应使用专门算法。

容量为 long long，距离及费用为 int128；费用不能为 LLONG\_MIN，因反边须取相反数。
另要求总流不超过 long long，且 $\sum cap\,|cost|<2^{120}$。
P3381 另保证答案不超过 $2^{31}-1$，其示例才可转换为 long long 输出。
SPFA 版实际增广 $A$ 次的复杂度为 $O((A+1)VE)$，不能假设 $A=O(V)$；
通常优先势能 Dijkstra，本版保留为直接的备选写法。

两版通过 $4000$ 组独立流量枚举及边界测试，完整驱动通过最大规模本地检查。
验证细节见仓库 docs/FLOW-ATTACHMENT.md；本地通过不等于线上 AC 或速度排名。

\section{2-SAT 编号方向与双连通定义}\label{knowledge-twosat-bcc}
TwoSAT 的 add(x,a,y,b) 表示 $(x=a)\lor(y=b)$。
强制值写 add(x,a,x,a)，蕴含 $(x=a)\Rightarrow(y=b)$ 写 add(x,!a,y,b)。
本库使用 Kosaraju 拓扑编号，真文字所在分量编号更大时取真；
附件 Tarjan 的分量编号为逆拓扑序，比较方向相反。
不能更换 SCC 实现而不检查赋值的不等号。
solve 可在追加子句后重算，仅成功时读取 ans；一般复杂度为 $O(n+m)$，不保证字典序最小。

点双采用连通且无割点的定义：桥两端也成块，孤立点单独成块，割点允许属于多个块。
自环不改变顶点间连通性，不额外产生点双；原边号仍保留，且不标桥。
每条逻辑无向边只调用一次 add，不能把两个方向分别当作两条边加入。
共享 lowlink 只跳过父边编号，不能跳过所有指向父顶点的重边，否则会误判桥。
附件点双单独跳过父顶点虽不改变点双阈值的真假，但它的 low 不可直接复用来判桥。

BiconnectedCore 的点号为 $1\ldots n$、边号从零开始。
run 后 blocks 是点双顶点集，bridge 是原边桥标记，bel 是 $1\ldots cnt$ 的边双号；
边双编号不承诺拓扑方向。不连通输入的圆方图和桥缩点结果均为森林。
st 仅保存未归块顶点，不是模拟 DFS 的执行栈；每次 run 重新计算所有结果。

独立子集、删点、删边、真值表与最大规模驱动检查均有普通及 ASan/UBSan 记录。
LC 点双 $22$ 组、边双 $21$ 组官方数据双模式本地 checker 接受，不等于线上 AC。
百万变量矛盾链在 ASan 下需更大栈，本地以 $1$ GiB 线程栈运行原递归算法；
提交仍用普通 main。细节见仓库 docs/GRAPH-ATTACHMENT.md。
当前共享核心的边双输出已由 CCF\_NOI 提交 QOJ999 获 AC（记录3061643，87ms，42352kb）；
本次在线结果不核验 blocks、cut、森林或补边方案。

\Needspace{13\baselineskip}
SCC 的完整拓扑序输出见第~\pageref{usage-example-161} 页，注意其编号方向与 TarjanSCC 相反。
固定 LC scc 的 $12$ 组官方数据由本地 checker 接受；枚举三点有向图、独立传递闭包
以及 $50$ 万点长链、环、分离二元环核对分组和跨组顺序。
同一 Kosaraju 驱动已由 CCF\_NOI 提交 QOJ906 获 AC（记录3061625，226ms，76512kb）；
dag 接口及其他未调用的捆绑算法不在本次覆盖范围，速度排名尚未核验。

已有点双分解时，removal\_components 返回原图块数 before 和独立删点后的总块数 after。
删除孤立点会使总数减一；非单点点双每包含该顶点一次，给删除后的所在连通分量多贡献一块。


以负块数和原点号排序，别把输入的输出名额 m 当作边数，保留每组输出后的空行。
独立逐点删除图并重新求连通块，检查数量及同分顺序。

连通无向图能强连通定向当且仅当无桥。
必要性：桥不论朝哪边定向，都不能双向通过；充分性：DFS 树边向下，非树边向祖先，
每个非根子树都有返回严格祖先的边，逐层可回到根，根又能沿树边到达所有顶点。

完整驱动枚举至五点的 $771$ 个连通简单图，以逐边删除判可行性，再检查每条输出边及正反可达性；
$10$ 万点、$30$ 万边边界和重边、自环的额外契约测试均通过。
上述测试均有普通及 ASan/UBSan 记录，算法保留递归，不是线上 AC 或速度排名。


\Needspace{15\baselineskip}
圆方森林保留全部原点 $1\ldots n$，第 $i$ 个点双对应方点 $n+i+1$（$i$ 从零开始）。

对于同一连通分量中的不同原点 $s,f$，圆方树路径经过的各块，恰好给出可能作为中间点的顶点并集。
同一块内任意第三点可放在两个指定端点间的简单路径上；块间只通过割点连接。
方点权先计块大小，路径内部割点的 $-1$ 去重，两端的 $-1$ 排除 $s,f$，
故路径权和就是可选中间点数，而不是简单路径条数。

设所在树有 $C$ 个原点，删去节点 $u$ 后各支的原点数为 $a_i$。
有序原点端点对经过 $u$ 的次数为
$2\sum_{i<j}a_i a_j+2(C-1)[u\le n]$；后项计端点等于 $u$ 的情形。
sz 仅统计原点，父侧大小为 $C-sz[u]$，遍历子树时用前缀和累计两两乘积。
乘以节点权后求和，使用 long long。每棵树分别求 $C$，孤立点贡献零；总时间空间 $O(n+m)$。

桥森林邻接项是（边双号，原桥号），不是原点端点；原点所属边双读 bel。

连通且允许新增重边时，每个叶子至少需一个新增边端点，故至少 $\lceil L/2\rceil$ 条边。
DFS 顺序收集叶子，奇数时补首叶，再前后半段配对；每条桥两侧的叶子区间都有配对跨过，故达到下界。
仅一个边双时不加边。非连通森林不能直接套此公式或调用构造接口。
P2860 只检查条数，方案端点及桥编号另以本地逐边删除和多重集检验；小图额外枚举全部新增边多重集验证最优性。
P4630、P2860 都是比赛应用，不代替正式模板题覆盖或线上 AC。

